We give a brief summary of some of the conceptual steps that led to today's mathematical understanding of the universal approximation theorems. For a thorough and extensive historical bibliography on neural networks, \cite{schmidhuber2015deep} compiled 888 references. See \cite{goodfellow2016deep} for an elementary well documented introduction to the subject.

\section{Universal approximation theorems}

A neural network (NN; for a definition, see e.g. \cite{pinkus1999approximation} or \cite{calin2020deep}) transforms an input $x\in D\subset\mathbb{R}^{n}$ into an output $\widetilde{y}=\Phi\left(  x\right)  \in\mathbb{R}$. As such it is a concrete real function of $n$ real variables, $n$ being a priori an arbitrary positive integer. A natural question that has been asked for a long time is the following: can this NN (with an appropriate architecture) approximate functions of a given class (common example: multivariate continuous functions defined on $D=K$, a compact subset of $\mathbb{R}^{n}$, but more general function spaces, such as $L^{p}$ spaces, are usually the proper framework from an abstract mathematical viewpoint) \citep{mhaskar1996system}? In other words, can a NN generate a topologically dense subset of the function class in question? For example, in the case of a normed vector space, given $\epsilon>0$, is there and, if so, can one actually construct at a reasonable cost a NN such that

\[
\left\Vert f-\widetilde{y}\right\Vert <\epsilon\text{ ?}%
\]

A related question (among many other researchers are addressing \citep{fan2020universal}) is: to what extent does the choice of the norm (topology) affect this existence and/or construction?

The proper mathematical interpretation of NN is based on the universal
approximation theorems (UAT). The UAT are roughly divided into two categories:

\begin{enumerate}
\item arbitrary number of neurons in the case of {\it shallow} neural networks (SNN, arbitrary width). This is the classical case. It goes back to the pioneering work of the 1940-1980's;

\item arbitrary number of layers with {\it deep} neural networks (DNN, arbitrary
depth). In practice, the latter have proven to be more efficient
\citep{mhaskar2019function}. It is always possible to build a SNN to emulate a DNN. However the number of neurons grows exponentially in order to obtain a
comparable performance (in terms of error) \citep{telgarsky2016benefits,eldan2016power}.
\end{enumerate}

In this sense, the UAT are representations of classical existence theorems. They implement the construction (i.e. find the appropriate weights) of the approximation of large classes of functions \citep{paluzo2020two}. However, they do not say how one explicitly (with a formula or an algorithm) finds this approximation (weights).

We give a brief summary of some of the conceptual steps that led to today's mathematical understanding of the UAT.

\begin{itemize}
\item The Kolmogorov--Arnold representation (or superposition) theorem \citep{kolmogorov1957representation} states that any continuous real multivariate function can be represented by a superposition of real one-variable continuous functions $f:\mathbb{R}^{n}\rightarrow\mathbb{R}$:
\[
f(x)  =\sum_{i=0}^{2n}\Phi_{i}\left(  \sum_{j=1}^{n}\phi
_{i,j}(x_{j})  \right)
\]
This theorem gives a partial answer to Hilbert's 13th problem \citep{akashi2001application,khesin2014arnold}. A construction is proposed by \cite{polar2020urysohn}.

\item Prior to the 1980s \citep{mcculloch1943logical,fukushima1975cognitron}, NN architectures were shallow. An understanding of SNN gives a good appreciation of DNN.

\item The Kolmogorov-Arnold theorem gives an exact representation of the function $f(x)$. This representation proves to be very complex (fractal nature) and computationally very slow, therefore unusable in NN. The accuracy condition had to be relaxed for an approximation condition ($L^{p}$
or uniform norm). The optimal solution is sought by iterative methods (the
gradient descent, introduced by Jacques Hadamard in 1908, being the most
popular). Pioneering modern work is due to \cite{cybenko1989approximation} for
sigmoid activation functions (arbitrary width), \cite{funahashi1989approximate}, \cite{hornik1991approximation} (he showed that UAT was
potentially constructible using a multilayer architecture of the NN rather
than a specific choice of the activation function), \cite{barron1993universal} (bounds on smooth function errors), etc...

\item  \cite{leshno1993multilayer} and \cite{pinkus1999approximation} have
shown that the universal approximation property (see \cite{kratsios2019universal} for a
definition) is equivalent to a non polynomial activation function.
\end{itemize}

\begin{theorem}
(\cite{cybenko1989approximation,hornik1991approximation,pinkus1999approximation}). Let $n\in\mathbb{N}$, $K\subset\mathbb{R}^{n}$ compact, $f:K\rightarrow\mathbb{R}$ continuous, $\rho:\mathbb{R}\rightarrow\mathbb{R}$ non polynomial and continuous, and let $\epsilon>0$. Then there exists a
two-layer NN $\Phi$ with activation function $\rho$ such that
\[
\left\Vert \Phi-f\right\Vert _{\infty}<\epsilon.
\]
\end{theorem}

We also note that DNN have demonstrated their capabilities in the approximation of functions \citep{liang2016deep,yarotsky2017error,fan2020universal}. They are found to have good experimental (empirical) performance, but it is largely unknown why. This is a very active area of research \cite{caterini2018deep}.

\begin{theorem}
(\cite{lu2017expressive}) Let  $f:\mathbb{R}^{n}\rightarrow\mathbb{R}$ be Lebesgue-integrable and $\epsilon>0$. Then there exists a ReLU
network of width $\leq n+4$ such that
\[
\int_{\mathbb{R}^{n}}\left\vert \Phi\left(  x\right)  -f\left(  x\right)  \right\vert dx<\epsilon
\]

\end{theorem}

\begin{definition}
(\cite{kidger2020universal}) Let $\rho:\mathbb{R}\rightarrow\mathbb{R}$ and $n,m,k\in\mathbb{N}$. Then let $NN_{n,m,k}$ represent the class of functions $\mathbb{R}^{n}\rightarrow\mathbb{R}^{m}$ described by feedforward neural networks with $n$ neurons in the input layer, $m$ neurons in the output layer, and an arbitrary number of hidden layers, each with $k$ neurons with activation function $\rho$. Every neuron in the output layer has the identity activation function.
\end{definition}

\begin{theorem}
(\cite{kidger2020universal}) Let $\rho:\mathbb{R}\rightarrow\mathbb{R}$ be any non-affine continuous function which is continuously differentiable at
at least one point, with non-zero derivative at that point. Let $K\subset\mathbb{R}^{n}$ be compact. Then NN$_{n,m,n+m+2}$ is dense in $C\left(  K,\mathbb{R}^{m}\right)$ with respect to the uniform norm.
\end{theorem}


\section{Low-rank approximation}

Due to its fundamental role in Machine Learning, a huge literature is devoted
to the singular value decomposition (SVD) of a matrix. We refer to some of the
well respected texts: \cite{strang1993introduction}, \cite{trefethen1997numerical}, \cite{watkins2004fundamentals}, \cite{gentle2007matrix}, \cite{elden2007matrix}, \cite{lange2010numerical},
and \cite{yanai2011projection} for
instance. Let $A\in \mathbb{R}^{m\times n}$ be a matrix of rank $r$. Its SVD consists in factoring $A$ as
\[
A=U\Sigma V^\top=\sum_{i=1}^{r}\sigma_{i}\u_{i}\v_{i}^\top
\]
where $U\in\mathbb{R}^{m\times m}$ and $V\in\mathbb{R}^{n\times n}$ are orthogonal square matrices, and $\Sigma\in\mathbb{R}^{m\times n}$ (\textquotedblleft diagonal\textquotedblright\ rectangular matrix of the ordered singular values: $\sigma_{1}\geq\ldots\geq\sigma _{r}>\sigma_{r+1}=\ldots=0=\sigma_{\min\{m,n\}}$). This
decomposition, and the related principal component analysis (PCA), has several
uses in Data Analysis. Written as $AV=U\Sigma$, SVD says that there exist orthogonal matrices $U$ and $V$ such that $A$ maps the columns $\v_{i}$ of $V$ (right singular vectors of $A$, directions or principal axes in the PCA parlance) into the columns $\u_{i}$ of $U$ (left singular vectors of $A$) such as
\[
A\v_{i}=\sigma_{i}\u_{i},
\]
the principal components.

Given a positive integer $1\leq k\leq r$, introduce the truncated matrices
$U_{k}=\left[  \u_{1},\ldots,\u_{k}\right]  \in \mathbb{R}^{m\times k}$, $V_{k}=\left[  \v_{1},\ldots,\v_{k}\right]  \in \mathbb{R}^{k\times n}$, and $\Sigma_{k}=\mathrm{diag}(  \sigma_{1},\ldots,\sigma_{k}) \in \mathbb{R}^{k\times k}$. Then the matrix of rank $k$
\[
A_{k}=U_{k}\Sigma_{k}V_{k}^\top=\sum_{i=1}^{k}\sigma_{i}\u_{i}\v_{i}^\top
\]
is an optimal lower-rank approximation of $A$ in the sense of both the
Frobenius norm and the $L_2$-norm. Indeed, the Eckart--Young--Mirsky theorem (for
a proof, see \cite{markovsky2012low}) states that among all matrices $B\in\mathbb{R}^{m\times n}$ of a given rank $k\leq r$, the one closest to $A$ is
$\hat{A}=A_{k}$, i.e.

\begin{align*}
\inf_{\mathrm{rank}(B)=k}\left\Vert A-B\right\Vert _{F} &  =\left\Vert
A-\hat{A}\right\Vert _{F}=\sqrt{\sum_{i=k+1}^{r}\sigma_{i}^{2}},\\
\inf_{\mathrm{rank}(B)=k}\left\Vert A-B\right\Vert _{2} &  =\left\Vert
A-\hat{A}\right\Vert _{2}=\sigma_{k+1}.
\end{align*}

Furthermore, the approximating matrix $\hat{A}$ is unique if and only if
$\sigma_{k+1}\neq\sigma_{k}$ \citep{markovsky2012low}. The theorem was proved in 1907 by Schmidt (of the Gram-Schmidt procedure) for the Frobenius norm \citep{leon2013gram}. A constructive solution was given by \cite{eckart1936approximation}. This solution was
shown by \cite{mirsky1960symmetric} to hold for all unitarily invariant norms. A further generalization addressing collinearity issues was proposed by \cite{golub1987generalization}. See \cite{stewart1993early} for historical details. 

One interpretation of the decomposition $A_{k}=U_{k}\Sigma_{k}V_{k}^{T}$ is as
a compression/learning linear procedure of the data matrix $A$ (in other
words, a dimension reduction method). Indeed, since the SVD extracts data in
the direction with the highest variance (and going down to those with the
lowest variance), matrix $U_{k}$ is the compression matrix in the first $k$
directions, matrix $\Sigma_{k}$ contains the amount of error (cost function)
we tolerate, and $V_{k}$ is the decompression matrix. If we keep all the non
zero singular values ($k=r$), there is theoretically no loss in the
compression (assuming that the data in $A$ are exact). If we drop the small
singular values (set them to $0$), it is not possible to reconstruct the
initial matrix $A$ exactly, but we obtain a matrix $A_{k}$ that is close. In
other words, the SVD decomposition is a way to control the loss of information
from $A$ under a predetermined level.

The assumption (routinely validated in image processing, for example) in the
above approach is that the components (columns of $U$) of the small singular
values contain little information (mostly \textquotedblleft
noise\textquotedblright) and should be discarded. However, the low-rank matrix approximation is easily seen not to be optimal if used to solve the linear regression problem (hence as a predictive tool) $y=A\beta+\epsilon$, that is to minimize
\begin{equation}
    \inf_{\mathrm{rank}(B)=k}\left\Vert \y-B\bbeta\right\Vert . 
    \label{eq:svd_regression}
\end{equation}

While the SVD picks the parameter $k-$vector $\bbeta$ of minimum norm
\citep{elden2007matrix,lange2010numerical},

\begin{equation}
    \hat{\bbeta}=A_{k}^{\dag} \y=\sum_{i=1}^{k}\frac{1}{\sigma_{i}} \u_{i}^\top \y \v_{i},
    \label{eq:svd_reg_solution}
\end{equation}

where $A_{k}^{\dag}$ is the pseudoinverse of $A_{k}$, the general solution of
\eqref{eq:svd_regression} is unknown \citep{xiang2012optimal}. See also \cite{talwalkar2010matrix} for a discussion in the Machine Learning context.

Some authors \citep{jolliffe1982note,hadi1998some} warned against the misconceived (mainly due to collinearity considerations) use of \eqref{eq:svd_reg_solution} in an ordinary linear regression setting. Dropping some of the small singular values may result in ignoring significant variables, mainly because the approximation does not take into account the response variable $\y$. A partial solution is provided by the ridge regression \citep{hoerl1970ridge}. However, still from a predictive viewpoint, SVD may prove to be an interesting approach in deep neural networks \citep{bermeitinger2019singular}. Used for initialization in the first layer, SVD gave promising experimental results. Theoretical issues still need to be addressed.

\newpage